\chapter{Extension Runtime}
\label{ch:extension}

The gert v2 extension runtime defines how external packages contribute tools,
schema fields, governance policies, and input providers to a running gert host
without requiring changes to the host binary.  Extensions run as child
processes; they communicate with the host over JSON-RPC 2.0 and are governed
by an explicit capability grant model.

A normative reference specification is maintained at
\texttt{specs/002-extension-runtime-v0/spec.md} in the repository.  This
chapter derives from and supersedes that document for design purposes; in any
conflict between the two, the specification takes precedence for
implementation.

\section{Architecture Overview}

\begin{enumerate}
  \item The gert host discovers extension manifests on startup.
  \item The host evaluates compatibility and capability policy before launching
        any process.
  \item Each accepted extension is started as a child process; its executable
        communicates with the host over its declared transport.
  \item The host and extension perform a versioned handshake.  The host sends
        the set of capabilities it is willing to grant; the extension confirms
        receipt.
  \item The extension registers its contributions (tools, schema fields,
        policies, providers).
  \item For the remainder of the run the host dispatches invocations to the
        extension and enforces all capability constraints.
  \item At run completion the host sends a graceful shutdown signal.
\end{enumerate}

\section{Capability Taxonomy}
\label{sec:ext-capabilities}

Every operation an extension may perform is gated by a named capability.
Capabilities are requested in the manifest and granted (or denied) by host
policy before the process is started.  The initial capability set is given in
Table~\ref{tab:capabilities}.

\begin{table}[ht]
\centering
\begin{tabular}{p{0.32\linewidth}p{0.30\linewidth}p{0.30\linewidth}}
\hline
\textbf{Capability name} & \textbf{Allows} & \textbf{Excludes} \\
\hline
\texttt{capability/tool-registration} &
  Register new tool definitions at runtime via \texttt{contributions/list}. &
  Does not allow the extension to invoke other tools or modify existing tool
  definitions. \\
\hline
\texttt{capability/schema-extension} &
  Contribute namespaced additional fields to the runbook schema (e.g.\
  \texttt{x-acme.*} annotations). &
  Cannot modify or remove existing schema fields; cannot alter core validation
  rules. \\
\hline
\texttt{capability/event-subscription} &
  Subscribe to host runtime events (step-started, step-completed, run-ended,
  etc.) via the \texttt{events/subscribe} method. &
  Cannot emit synthetic events into the host event bus. \\
\hline
\texttt{capability/policy-contribution} &
  Register additional governance policy rules evaluated during pre-flight and
  approval checks. &
  Cannot override or disable built-in governance rules; cannot grant
  capabilities to other extensions. \\
\hline
\texttt{capability/provider-registration} &
  Register new input providers that handle \texttt{from:} prefix bindings. &
  Cannot intercept resolution requests destined for another registered
  provider. \\
\hline
\texttt{capability/file-read} &
  Read files within the path scope declared in the manifest
  (\texttt{file-read-paths}). &
  Cannot read files outside declared paths; cannot write or delete files. \\
\hline
\texttt{capability/file-write} &
  Write or create files within the path scope declared in the manifest
  (\texttt{file-write-paths}). &
  Cannot write outside declared paths; does not imply read access. \\
\hline
\texttt{capability/network} &
  Make outbound network calls to the hosts declared in the manifest
  (\texttt{network-hosts}). &
  Cannot connect to undeclared hosts; cannot listen for inbound connections. \\
\hline
\texttt{capability/env-read} &
  Read environment variables whose names match the declared
  \texttt{env-read-patterns} list. &
  Cannot read undeclared variables; cannot write or unset environment
  variables. \\
\hline
\texttt{capability/run-state-read} &
  Read current run state (step status, captured variables, trace events). &
  Cannot modify run state; does not imply the ability to emit output. \\
\hline
\end{tabular}
\caption{Extension capability taxonomy (v2).}
\label{tab:capabilities}
\end{table}

The capability string format is \texttt{capability/<name>}.  Future versions
may introduce scoped variants using path notation, e.g.\
\texttt{capability/file-read:logs/**}.

\section{Extension Manifest Format}
\label{sec:ext-manifest}

Each extension ships a \texttt{gert-extension.yaml} file alongside its
executable.  The manifest declares identity, requested capabilities, transport,
and compatibility constraints.

\begin{verbatim}
# gert-extension.yaml — Extension Manifest
apiVersion: extension/v2

meta:
  name: acme.incident-pack          # Reverse-DNS namespaced identifier
  version: "1.2.0"                  # Semver
  description: "Acme incident resolution tools and providers"
  author: "Acme Platform Team <platform@acme.example>"

compatibility:
  api-version: ">=2.0.0 <3.0.0"    # Host API version range (semver)
  protocol-version: "2.0"           # Extension protocol version

entry-point:
  executable: "./bin/acme-extension"
  args: ["--config", "acme.yaml"]

transport: stdio-jsonrpc             # stdio-jsonrpc | grpc | mcp

capabilities:
  - capability/tool-registration
  - capability/provider-registration
  - capability/schema-extension
  - capability/network
  - capability/file-read

# Required when capability/file-read is declared:
file-read-paths:
  - "./runbooks/**"
  - "./data/**"

# Required when capability/network is declared:
network-hosts:
  - "api.pagerduty.com"
  - "*.acme.example"

# Required when capability/env-read is declared:
env-read-patterns:
  - "ACME_API_*"
  - "PD_TOKEN"

platform:                            # Optional OS/arch constraints
  os: ["linux", "darwin", "windows"]
  arch: ["amd64", "arm64"]
\end{verbatim}

Manifest validation is performed by the host before the process is started.
An extension with an unsatisfied \texttt{compatibility.api-version} range is
rejected before any subprocess is created.

\section{Extension Discovery}
\label{sec:ext-discovery}

The host discovers extensions using the following resolution order, with later
entries taking precedence:

\begin{enumerate}
  \item \textbf{Built-in extensions}: compiled into the host binary; always
        available; no manifest required.
  \item \textbf{Workspace config}: extensions declared in
        \texttt{.gert/extensions.yaml} at the workspace root.
  \item \textbf{Runbook declaration}: extensions listed in the runbook's
        top-level \texttt{extensions:} field.
  \item \textbf{CLI override}: \texttt{--extension <path>} flags provided at
        invocation time.
\end{enumerate}

The \texttt{.gert/extensions.yaml} workspace config format:

\begin{verbatim}
# .gert/extensions.yaml
extensions:
  - path: "./extensions/acme-pack"   # directory containing gert-extension.yaml
  - path: "/opt/gert-extensions/audit-pack"
\end{verbatim}

The runbook \texttt{extensions:} field:

\begin{verbatim}
apiVersion: runbook/v2
meta:
  name: incident-triage
extensions:
  - path: "./extensions/acme-pack"
\end{verbatim}

When the same extension appears in multiple discovery sources, the host
deduplicates by \texttt{meta.name}.  If two distinct extensions share the same
name, the one with the higher resolution precedence wins; the lower one is
logged as a warning and skipped.

\section{Handshake and Lifecycle Protocol}
\label{sec:ext-handshake}

All messages use JSON-RPC 2.0.  The host is always the initiator; the
extension never sends unsolicited requests.

\subsection{extension/initialize (host \textrightarrow{} extension)}

Sent immediately after the extension process starts.  The host communicates
the capabilities it is willing to grant and its own version identity.

Request params:

\begin{verbatim}
{
  "jsonrpc": "2.0",
  "id": "1",
  "method": "extension/initialize",
  "params": {
    "host": {
      "name": "gert",
      "version": "2.0.0",
      "apiVersion": "2.0"
    },
    "protocolVersion": "2.0",
    "grantedCapabilities": [
      "capability/tool-registration",
      "capability/network"
    ],
    "runContext": {
      "runId": "run-abc123",
      "mode": "real"
    }
  }
}
\end{verbatim}

Success response:

\begin{verbatim}
{
  "jsonrpc": "2.0",
  "id": "1",
  "result": {
    "extension": {
      "name": "acme.incident-pack",
      "version": "1.2.0"
    },
    "protocolVersion": "2.0",
    "acknowledgedCapabilities": [
      "capability/tool-registration",
      "capability/network"
    ]
  }
}
\end{verbatim}

Error codes for this method:

\begin{itemize}
  \item \texttt{-32001}: Protocol version not supported by extension.
  \item \texttt{-32002}: Required capability not granted (extension considers
        the capability non-negotiable).
  \item \texttt{-32003}: Host API version outside extension compatibility range.
\end{itemize}

\subsection{contributions/list (host \textrightarrow{} extension)}

Sent after a successful \texttt{extension/initialized} notification.  The
extension returns its contribution manifest, filtered to the granted capability
set.

\begin{verbatim}
{
  "jsonrpc": "2.0",
  "id": "2",
  "method": "contributions/list",
  "params": {}
}
\end{verbatim}

Response:

\begin{verbatim}
{
  "jsonrpc": "2.0",
  "id": "2",
  "result": {
    "tools": [ { "name": "acme.pd-lookup", ... } ],
    "providers": [ { "name": "pd", "prefixes": ["pd."] } ],
    "schemaExtensions": [],
    "policyContributions": []
  }
}
\end{verbatim}

\subsection{extension/ping (host \textrightarrow{} extension)}

Periodic health check.  The extension MUST respond within the host-configured
ping timeout (default 5 seconds).  A non-response triggers a crash-recovery
cycle (Section~\ref{sec:ext-crash}).

\begin{verbatim}
{ "jsonrpc": "2.0", "id": "42", "method": "extension/ping", "params": {} }
\end{verbatim}

Response:

\begin{verbatim}
{ "jsonrpc": "2.0", "id": "42", "result": { "status": "ok" } }
\end{verbatim}

\subsection{extension/shutdown (host \textrightarrow{} extension)}

Graceful termination.  The extension SHOULD drain in-flight requests and
release resources within the shutdown timeout (default 10 seconds).

\begin{verbatim}
{ "jsonrpc": "2.0", "id": "99", "method": "extension/shutdown", "params": {} }
\end{verbatim}

Response (best-effort; host may not wait):

\begin{verbatim}
{ "jsonrpc": "2.0", "id": "99", "result": { "status": "ok" } }
\end{verbatim}

If the extension does not exit within the shutdown timeout, the host sends
SIGTERM followed by SIGKILL after a grace period.

\subsection{Lifecycle State Machine}

\begin{enumerate}
  \item \textbf{discovered}: manifest found and parsed; no process yet.
  \item \textbf{verified}: manifest validated; capabilities evaluated against
        policy.
  \item \textbf{starting}: executable launched; stdin/stdout pipes attached.
  \item \textbf{initializing}: \texttt{extension/initialize} sent; awaiting
        response.
  \item \textbf{ready}: \texttt{contributions/list} completed; accepting
        invocations.
  \item \textbf{draining}: \texttt{extension/shutdown} sent; waiting for
        in-flight completions.
  \item \textbf{stopped}: process has exited cleanly.
  \item \textbf{crashed}: process exited unexpectedly (see
        Section~\ref{sec:ext-crash}).
\end{enumerate}

\section{Versioning and Compatibility}
\label{sec:ext-versioning}

The host maintains a monotonically increasing API version string following
semver.  The extension manifest declares a semver range in
\texttt{compatibility.api-version}.  Compatibility is evaluated before any
process is started:

\begin{itemize}
  \item If the host API version falls outside the declared range, the extension
        is rejected with a structured error logged to the run trace.  The run
        continues with the remaining available extensions; a warning is surfaced
        to the operator.
  \item During \texttt{extension/initialize}, the host also sends its
        \texttt{protocolVersion}.  If the extension does not support that
        protocol version it MUST return error code \texttt{-32001}.
  \item If an extension requests a capability that the host version does not
        implement (e.g., a capability introduced in a later host version), the
        capability is silently absent from \texttt{grantedCapabilities}.  The
        extension SHOULD treat absent non-required capabilities as optional
        degradation rather than a fatal error.
\end{itemize}

Host API versions are tracked in \texttt{specs/api-versions.md}.  Each new
capability addition constitutes a minor version bump; breaking changes to
existing method contracts constitute a major version bump.

\section{Sandboxing and Process Isolation}
\label{sec:ext-sandbox}

\subsection{Process Isolation}

Extensions always run as child processes of the gert host.  They are never
loaded as shared libraries or plug-ins executed in-process.  The host:

\begin{itemize}
  \item Sets a restricted environment: only env vars matching
        \texttt{env-read-patterns} from the manifest are forwarded.
  \item Attaches the extension's stderr to a structured log collector; raw
        stderr lines are stored in the run trace under
        \texttt{extension.<name>.stderr}.
  \item Does not forward the host process's file descriptor table beyond
        stdin/stdout/stderr.
\end{itemize}

\subsection{Capability Enforcement}

Every RPC method the extension may call is guarded by a capability check.  The
host maintains a granted-capability set per extension process.  Before
dispatching any extension-initiated request (e.g., a tool registration during
\texttt{contributions/list}) the host verifies the required capability is in
the granted set.  Violations return:

\begin{verbatim}
{
  "jsonrpc": "2.0",
  "id": "...",
  "error": {
    "code": -32010,
    "message": "capability not granted",
    "data": {
      "required": "capability/file-write",
      "granted": ["capability/file-read", "capability/network"]
    }
  }
}
\end{verbatim}

File-path and network-host constraints are enforced by the host at the system
call boundary using OS-level mechanisms where available (e.g., pledge/unveil on
OpenBSD, seccomp on Linux) and by host-side validation of all path and URL
arguments on other platforms.

\subsection{Timeout and Crash Handling}
\label{sec:ext-crash}

\begin{itemize}
  \item \textbf{Per-invocation timeout}: each \texttt{tools/invoke} call
        carries a \texttt{timeoutMs} field.  If the extension does not respond
        within that deadline, the host sends a \texttt{tools/cancel} request
        and, if still unresponsive after a grace period, terminates the process.
  \item \textbf{Ping timeout}: if an extension fails to respond to
        \texttt{extension/ping} within 5 seconds it is considered unresponsive.
        The host logs a structured warning and terminates the process.
  \item \textbf{Unexpected exit}: if the extension process exits with a
        non-zero status or is killed by a signal, the host marks the extension
        as \textbf{crashed}, emits a \texttt{extension.crashed} trace event,
        and cancels all in-flight invocations with a structured
        \texttt{-32099} (extension crashed) error.
  \item \textbf{Host isolation}: extension crashes never crash the host.  A
        run that loses an extension mid-execution surfaces an error on the
        affected step and allows the operator to decide whether to abort or
        continue with a degraded tool set.
\end{itemize}
